\scanpage{112}
\mspage{57}
% archive sheet number: 111

\noindent\begin{minipage}{.62\linewidth}\small
\diagram{\begin{tikzcd}[ampersand replacement=\&,row sep=large]
 \underline{\mathrm{Parf}}\,S\arrow[d,"r"']\\
 \mathrm{Cat}_{/S}\arrow[u,bend right=38,dashed,"C"']
\end{tikzcd}}
\end{minipage}\hfill\begin{minipage}{.36\linewidth}\small
\[
 \left\{\begin{aligned}rC&\longleftarrow\mathrm{id}\\ Cr&\longleftarrow\mathrm{id}\end{aligned}\right.
\]
\end{minipage}\par

\noindent\begin{minipage}{.62\linewidth}\small
\diagram{\begin{tikzcd}[ampersand replacement=\&,column sep=large,row sep=large]
 \underline{\mathrm{Fibsc}}\,S\arrow[dr,"\lambda"]\&\&\underline{\mathrm{Cofibsc}}\,S\arrow[dl,"\mu"']\\
 \&\mathrm{Cat}_{/S}\arrow[ul,bend left=45,dashed,"\Phi"']\arrow[ur,bend right=45,dashed,"\Psi"]\&
\end{tikzcd}}
\end{minipage}\hfill\begin{minipage}{.36\linewidth}\small
\[
 \begin{aligned}
 &\left\{\begin{aligned}\lambda\Phi&\longleftarrow\mathrm{id}\\\Phi\lambda&\longrightarrow\mathrm{id}\end{aligned}\right.\\[3pt]
 &\left\{\begin{aligned}\mu\Psi&\longleftarrow\mathrm{id}\\\Psi\mu&\mathbin{\unsure{\to|-}}\mathrm{id}\end{aligned}\right.
 \end{aligned}
\]
\end{minipage}\par

\noindent\begin{minipage}{.62\linewidth}\small
\diagram{\begin{tikzcd}[ampersand replacement=\&,column sep=large,row sep=large]
 \underline{\mathrm{Fibsc}}_0\,S\arrow[dr,"\lambda_0"]\&\&\underline{\mathrm{Cofibsc}}_0\,S\arrow[dl,"\mu_0"']\\
 \&\underline{\mathrm{Fibhot}}\,S\arrow[ul,bend left=45,dashed,"\Phi_0"']\arrow[ur,bend right=45,dashed,"\Psi_0"]\&
\end{tikzcd}}
\end{minipage}\hfill\begin{minipage}{.36\linewidth}\small
\[
 \begin{aligned}
 &\left\{\begin{aligned}\lambda_0\Phi_0&\longleftarrow\mathrm{id}\\\Phi_0\lambda_0&\longrightarrow\mathrm{id}\end{aligned}\right.\\[3pt]
 &\left\{\begin{aligned}\mu_0\Psi_0&\longleftarrow\mathrm{id}\\\Psi_0\mu_0&\mathbin{\unsure{-|\to}}\mathrm{id}\end{aligned}\right.
 \end{aligned}
\]
\end{minipage}\par

\margin{\uline{NB} Tous les foncteurs marqués sont compatibles avec la $W$-équivalence grossière, et les dix hom fonctoriels sont des $W$-équivalences grossières.}

NB les quatre derniers hom. de foncteurs, relatifs au troisième diagramme (avec les indices 0) sont simplement induits par ceux relatifs au diagramme qui le précède. Donc les vérifications-clef sont relatives aux \deleted{\unsure{deux}} premiers diagrammes, et par symétrie, on n'a finalement qu'à regarder les \uline{deux} couples de foncteurs $(r,C),(\lambda,\Phi)$, et pour \unsure{chacun|chacune} de ces couples, exhiber les deux hom fonctoriels, des deux \unsure{composés|comparés} avec les foncteurs identiques.

